\documentclass[12pt]{scrartcl}
\usepackage[hagavi]{azzam}

% Modulo dan Keterbagian - Modul (gabungan Materi dan Soal Latihan)
% Digabung dari Modulo-Keterbagian_Materi.tex dan Modulo-Keterbagian_Latihan.tex.
% Adaptasi dari 8 Outsource K Math/Pra-OSK/Keterbagian.tex dan Modulo.tex.
% _Ujian dan _Solusi tetap terpisah.
% Sumber tiap soal latihan dimunculkan di PDF sebagai \item (Sumber); komentar menyimpan detail lengkap.
% Soal pilihan ganda asli disajikan tanpa opsi jawaban sebagai soal isian,
% kecuali soal konsep nomor 1 yang memang membutuhkan opsinya.

\title{Modulo dan Keterbagian - Modul}
\date{15 Oktober 2026}

\begin{document}
\maketitle

\renewcommand*\contentsname{Daftar Isi}
\tableofcontents

\newpage

\part*{Materi}

Modul ini merangkum dua fondasi teori bilangan yang hampir selalu muncul di OSK dan OSP:
keterbagian dan kongruensi modulo. Setiap teorema langsung diikuti contoh soal,
jadi bacalah teori lalu contohnya berpasangan.
Kuasailah sifat dasarnya sampai hafal di luar kepala,
karena hampir setiap soal teori bilangan tingkat lanjut memakainya tanpa menyebutnya lagi.

\section{Keterbagian}

\begin{definition}
Diberikan dua bilangan bulat $a$ dan $b$ dengan $a \ne 0$.
Kita menyebut $a$ \textbf{membagi} $b$ apabila terdapat bilangan bulat $c$
sedemikian hingga $b = ca$. Notasinya adalah $a \mid b$.
\end{definition}

\begin{fact}[Sifat keterbagian]
Untuk bilangan bulat $a, b, c$ dengan pembagi tak nol, berlaku:
\begin{enumerate}
    \item $a \mid 0$ dan $a \mid a$.
    \item Jika $a \mid b$ dan $b \mid c$, maka $a \mid c$.
    \item Jika $a \mid b$ dan $a \mid c$, maka
    $a \mid \alpha b + \beta c$ untuk setiap bilangan bulat $\alpha, \beta$.
    \item Jika $a \mid b$ dengan $b \ne 0$, maka $|a| \le |b|$.
    \item Jika $a \mid b$ dan $b \mid a$, maka $|a| = |b|$.
    \item Jika $a \mid b$, maka $a \mid bc$ untuk setiap bilangan bulat $c$.
\end{enumerate}
\end{fact}

\begin{example}
Carilah semua bilangan bulat $n$ yang memenuhi $n + 1 \mid 100$.
\end{example}
\begin{proof}[Solusi]
Syaratnya berarti $n + 1$ adalah faktor dari $100 = 2^2 \cdot 5^2$.
Faktor positifnya ada $9$ buah, dan masing-masing versi negatifnya juga memenuhi:
\begin{align*}
n + 1 &\in \{\pm 1, \pm 2, \pm 4, \pm 5, \pm 10, \pm 20, \pm 25, \pm 50, \pm 100\}, \\
n &\in \{-101, -51, -26, -21, -11, -6, -5, -3, -2, 0, 1, 3, 4, 9, 19, 24, 49, 99\}.
\end{align*}
\end{proof}

\begin{example}
Carilah semua bilangan bulat $n$ yang memenuhi $n + 1 \mid n^2 + 99$.
\end{example}
\begin{proof}[Solusi]
Kurangkan kelipatan pembagi untuk menyisakan konstanta:
\[ n^2 + 99 = (n + 1)(n - 1) + 100. \]
Karena $n + 1$ selalu membagi $(n + 1)(n - 1)$, syaratnya setara dengan
\[ n + 1 \mid 100. \]
Jawabannya sama seperti contoh sebelumnya.
Trik menyisakan konstanta ini akan muncul berulang kali.
\end{proof}

\begin{theorem}[Algoritma Pembagian]
Untuk setiap bilangan asli $a$ dan $b$, terdapat pasangan \textbf{unik}
bilangan cacah $(q, r)$ yang memenuhi
\[ b = qa + r \quad \text{dengan} \quad 0 \le r < a. \]
\end{theorem}

\begin{example}
Buktikan bahwa di antara $n$ bilangan bulat berurutan, tepat satu habis dibagi $n$.
\end{example}
\begin{proof}[Solusi]
Ambil $n$ bilangan berurutan $m, m + 1, \dots, m + n - 1$ dan bagi $m$ dengan $n$:
\[ m = qn + r \quad \text{dengan} \quad 0 \le r < n. \]
Jika $r = 0$, maka $m$ sendiri yang habis dibagi $n$.
Jika $r > 0$, bilangan ke-$(n - r)$ setelah $m$ adalah
\[ m + (n - r) = qn + r + n - r = (q + 1)n, \]
yang habis dibagi $n$ dan masih di dalam barisan.
Tidak mungkin ada dua, karena selisih dua kelipatan $n$ yang berbeda minimal $n$,
sementara rentang barisan hanya $n - 1$.
\end{proof}

\subsection{Aturan keterbagian bilangan kecil}

\begin{enumerate}
    \item Sebuah bilangan asli habis dibagi $2$ apabila digit terakhirnya genap.
    \item Sebuah bilangan asli habis dibagi $4$ apabila dua digit terakhirnya habis dibagi $4$.
    \item Sebuah bilangan asli habis dibagi $8$ apabila tiga digit terakhirnya habis dibagi $8$.
    \item Secara umum, sebuah bilangan asli habis dibagi $2^n$ apabila
    $n$ digit terakhirnya habis dibagi $2^n$.
    \item Sebuah bilangan asli habis dibagi $5$ apabila digit terakhirnya $0$ atau $5$.
    \item Sebuah bilangan asli habis dibagi $25$ apabila dua digit terakhirnya
    $00$, $25$, $50$, atau $75$.
    \item Secara umum, sebuah bilangan asli habis dibagi $5^n$ apabila
    $n$ digit terakhirnya habis dibagi $5^n$.
    \item Sebuah bilangan asli habis dibagi $3$ maupun $9$ apabila jumlah digitnya
    habis dibagi $3$ maupun $9$.
    \item Sebuah bilangan asli habis dibagi $11$ apabila selisih jumlah digit
    di posisi ganjil dan jumlah digit di posisi genap habis dibagi $11$.
    \item Sebuah bilangan habis dibagi $10^n$ apabila $n$ digit terakhirnya nol.
    \item Sebuah bilangan habis dibagi $ab$ untuk $a$ dan $b$ yang relatif prima
    apabila bilangan tersebut habis dibagi $a$ dan juga habis dibagi $b$.
\end{enumerate}

Untuk bilangan $n$ dengan sedikitnya dua digit, tulis $f(n)$ untuk bilangan yang
diperoleh dengan menghapus digit terakhir $n$, dan $g(n)$ untuk digit terakhir $n$.
Sebagai contoh, $f(421) = 42$ dan $g(421) = 1$. Maka:
\begin{enumerate}
    \item $n$ habis dibagi $7$ apabila $f(n) - 2g(n)$ habis dibagi $7$.
    \item $n$ habis dibagi $13$ apabila $f(n) - 9g(n)$ habis dibagi $13$.
    \item $n$ habis dibagi $17$ apabila $f(n) - 5g(n)$ habis dibagi $17$.
    \item $n$ habis dibagi $19$ apabila $f(n) + 2g(n)$ habis dibagi $19$.
\end{enumerate}

\begin{example}
Buktikan bahwa setiap bilangan $6$ digit berbentuk $\overline{abcabc}$ habis dibagi $11$.
\end{example}
\begin{proof}[Solusi]
Uraikan bilangannya:
\begin{align*}
\overline{abcabc} &= 1000 \cdot \overline{abc} + \overline{abc} \\
&= 1001 \cdot \overline{abc} \\
&= 7 \cdot 11 \cdot 13 \cdot \overline{abc}.
\end{align*}
Jadi bilangan itu bahkan habis dibagi $7$ dan $13$ juga.
\end{proof}

\section{FPB dan KPK}

\begin{definition}
Untuk bilangan asli $m$ dan $n$, $FPB(m, n)$ adalah faktor persekutuan positif
terbesar keduanya, sedangkan $KPK(m, n)$ adalah kelipatan persekutuan positif
terkecil keduanya. Jika $FPB(m, n) = 1$, maka $m$ dan $n$ disebut relatif prima.
\end{definition}

\begin{fact}[Sifat FPB]
\begin{enumerate}
    \item Jika $d = FPB(m, n)$ dan $m = dm'$, $n = dn'$, maka $FPB(m', n') = 1$.
    \item Setiap faktor persekutuan dari $m$ dan $n$ membagi $FPB(m, n)$.
    \item Jika $m = nq + r$, maka $FPB(m, n) = FPB(n, r)$.
    \item Jika $m = p_1^{\alpha_1} \cdots p_k^{\alpha_k}$ dan
    $n = p_1^{\beta_1} \cdots p_k^{\beta_k}$, maka
    \[ FPB(m, n) = p_1^{\min(\alpha_1, \beta_1)} \cdots p_k^{\min(\alpha_k, \beta_k)}. \]
\end{enumerate}
\end{fact}

\begin{theorem}[Algoritma Euclid dan Bezout]
Dengan menerapkan algoritma pembagian berulang,
\[ m = nq_1 + r_1, \quad n = r_1q_2 + r_2, \quad r_1 = r_2q_3 + r_3, \quad \dots \]
sampai sisanya nol, sisa tak nol terakhir adalah $FPB(m, n)$.
Akibatnya, untuk setiap bilangan asli $m$ dan $n$ terdapat bilangan bulat $a$ dan $b$
sedemikian hingga $am + bn = FPB(m, n)$.
\end{theorem}

\begin{example}
Hitunglah $FPB(133, 198)$.
\end{example}
\begin{proof}[Solusi]
Terapkan algoritma Euclid:
\begin{align*}
198 &= 133 + 65, \\
133 &= 2 \cdot 65 + 3, \\
65 &= 21 \cdot 3 + 2, \\
3 &= 2 + 1.
\end{align*}
Sisa tak nol terakhir adalah $1$, jadi
\[ FPB(133, 198) = 1. \]
\end{proof}

\begin{example}
Tentukan bilangan bulat $a$ dan $b$ sehingga $133a + 198b = 1$.
\end{example}
\begin{proof}[Solusi]
Bentuk Bezout diperoleh dengan membaca langkah Euclid di atas mundur dari belakang:
\begin{align*}
1 &= 3 - 2 \\
&= 3 - (65 - 21 \cdot 3) \\
&= 22 \cdot 3 - 65 \\
&= 22 \cdot (133 - 2 \cdot 65) - 65 \\
&= 22 \cdot 133 - 45 \cdot 65 \\
&= 22 \cdot 133 - 45 \cdot (198 - 133) \\
&= 67 \cdot 133 - 45 \cdot 198.
\end{align*}
Jadi $a = 67$ dan $b = -45$ memenuhi.
\end{proof}

\begin{fact}[Sifat KPK]
\begin{enumerate}
    \item Setiap kelipatan persekutuan dari $m$ dan $n$ adalah kelipatan dari $KPK(m, n)$.
    \item Dengan notasi faktorisasi prima yang sama,
    \[ KPK(m, n) = p_1^{\max(\alpha_1, \beta_1)} \cdots p_k^{\max(\alpha_k, \beta_k)}. \]
    \item $FPB(m, n) \times KPK(m, n) = mn$.
\end{enumerate}
\end{fact}

\begin{example}
Hitunglah $KPK(24, 36)$ dan periksa $FPB(24, 36) \times KPK(24, 36) = 24 \cdot 36$.
\end{example}
\begin{proof}[Solusi]
Faktorkan kedua bilangan:
\begin{align*}
24 &= 2^3 \cdot 3, \\
36 &= 2^2 \cdot 3^2.
\end{align*}
Ambil pangkat terkecil untuk FPB dan terbesar untuk KPK:
\begin{align*}
FPB(24, 36) &= 2^2 \cdot 3 = 12, \\
KPK(24, 36) &= 2^3 \cdot 3^2 = 72.
\end{align*}
Pemeriksaan:
\[ 12 \cdot 72 = 864 = 24 \cdot 36. \]
\end{proof}

\section{Kongruensi Modulo}

\begin{definition}
Misalkan $a$, $b$, dan $n$ bilangan bulat dengan $n \ne 0$.
Kita menyebut $a$ \textbf{kongruen} dengan $b$ modulo $n$ apabila $n \mid a - b$.
Notasinya adalah $a \equiv b \pmod{n}$.
\end{definition}

\begin{fact}[Sifat kongruensi]
\begin{enumerate}
    \item $a \equiv a \pmod{n}$.
    \item Jika $a \equiv b \pmod{n}$, maka $b \equiv a \pmod{n}$.
    \item Jika $a \equiv b \pmod{n}$ dan $b \equiv c \pmod{n}$, maka $a \equiv c \pmod{n}$.
    \item Jika $a \equiv b \pmod{n}$, maka $ka \equiv kb \pmod{n}$ untuk setiap bilangan bulat $k$.
    \item Jika $a \equiv b \pmod{n}$ dan $c \equiv d \pmod{n}$, maka
    $ja + kc \equiv jb + kd \pmod{n}$ untuk setiap bilangan bulat $j$ dan $k$.
    \item Jika $a \equiv b \pmod{n}$ dan $c \equiv d \pmod{n}$, maka $ac \equiv bd \pmod{n}$.
    \item $a \equiv b \pmod{n_i}$ untuk $i = 1, \dots, k$ jika dan hanya jika
    $a \equiv b \pmod{KPK(n_1, \dots, n_k)}$.
\end{enumerate}
\end{fact}

\begin{example}
Tentukan sisa pembagian $2^{10} + 3^{10}$ oleh $7$.
\end{example}
\begin{proof}[Solusi]
Sederhanakan masing-masing suku modulo $7$:
\begin{align*}
2^3 = 8 &\equiv 1 \pmod{7}, \\
2^{10} = (2^3)^3 \cdot 2 &\equiv 1^3 \cdot 2 \equiv 2 \pmod{7}.
\end{align*}
\begin{align*}
3^2 = 9 &\equiv 2 \pmod{7}, \\
3^4 &\equiv 2^2 \equiv 4 \pmod{7}, \\
3^8 &\equiv 4^2 = 16 \equiv 2 \pmod{7}, \\
3^{10} = 3^8 \cdot 3^2 &\equiv 2 \cdot 2 \equiv 4 \pmod{7}.
\end{align*}
Jumlahkan:
\[ 2^{10} + 3^{10} \equiv 2 + 4 \equiv 6 \pmod{7}. \]
Sisanya $6$.
\end{proof}

\subsection{Teorema penting}

\begin{theorem}[Invers modular]
Untuk setiap bilangan asli $m$ dan $n$ yang relatif prima,
terdapat bilangan asli $x$ sedemikian hingga $mx \equiv 1 \pmod{n}$.
\end{theorem}

\begin{example}
Tentukan bilangan asli $x$ sehingga $7x \equiv 1 \pmod{31}$.
\end{example}
\begin{proof}[Solusi]
Terapkan Euclid pada $31$ dan $7$:
\begin{align*}
31 &= 4 \cdot 7 + 3, \\
7 &= 2 \cdot 3 + 1.
\end{align*}
Baca mundur untuk bentuk Bezout:
\begin{align*}
1 &= 7 - 2 \cdot 3 \\
&= 7 - 2 \cdot (31 - 4 \cdot 7) \\
&= 9 \cdot 7 - 2 \cdot 31.
\end{align*}
Modulo $31$ suku terakhir lenyap:
\[ 9 \cdot 7 \equiv 1 \pmod{31}. \]
Jadi $x = 9$ memenuhi (periksa: $63 = 2 \cdot 31 + 1$).
\end{proof}

\begin{theorem}[Euler]
Misalkan $\varphi(m)$ menyatakan banyak bilangan asli tidak lebih dari $m$
yang relatif prima dengan $m$. Jika faktorisasi prima dari $n$ adalah
\[ n = p_1^{\alpha_1} p_2^{\alpha_2} \cdots p_k^{\alpha_k}, \]
maka
\[ \varphi(n) = n\left(1 - \frac{1}{p_1}\right)\left(1 - \frac{1}{p_2}\right)
\cdots\left(1 - \frac{1}{p_k}\right). \]
Untuk $a$ dan $m$ yang relatif prima, berlaku $a^{\varphi(m)} \equiv 1 \pmod{m}$.
\end{theorem}

\begin{example}
Hitunglah $\varphi(100)$, lalu tentukan sisa $3^{41}$ dibagi $100$.
\end{example}
\begin{proof}[Solusi]
Faktorisasi $100 = 2^2 \cdot 5^2$ memberi
\[ \varphi(100) = 100\left(1 - \frac{1}{2}\right)\left(1 - \frac{1}{5}\right) = 100 \cdot \frac{1}{2} \cdot \frac{4}{5} = 40. \]
Karena $3$ relatif prima dengan $100$, teorema Euler memberi
\[ 3^{40} \equiv 1 \pmod{100}. \]
Maka
\[ 3^{41} = 3 \cdot 3^{40} \equiv 3 \cdot 1 \equiv 3 \pmod{100}. \]
Sisanya $3$.
\end{proof}

\begin{theorem}[Fermat Kecil]
Untuk setiap bilangan prima $p$ dan bilangan bulat $a$, berlaku
$a^p \equiv a \pmod{p}$. Jika $p \nmid a$, maka $a^{p-1} \equiv 1 \pmod{p}$.
\end{theorem}

\begin{example}
Carilah sisa pembagian $2^{100}$ oleh $13$.
\end{example}
\begin{proof}[Solusi]
Karena $13$ prima dan tidak membagi $2$, Fermat Kecil memberi
\[ 2^{12} \equiv 1 \pmod{13}. \]
Bagi pangkatnya:
\[ 100 = 12 \cdot 8 + 4. \]
Maka
\begin{align*}
2^{100} &= (2^{12})^8 \cdot 2^4 \\
&\equiv 1^8 \cdot 16 \\
&\equiv 3 \pmod{13}.
\end{align*}
Sisanya $3$.
\end{proof}

\begin{theorem}[Wilson]
Untuk setiap bilangan prima $p$, berlaku $(p-1)! \equiv -1 \pmod{p}$.
\end{theorem}

\begin{example}
Tentukan sisa pembagian $14!$ oleh $17$.
\end{example}
\begin{proof}[Solusi]
Wilson memberi
\[ 16! \equiv -1 \pmod{17}. \]
Uraikan $16!$:
\begin{align*}
16! &= 16 \cdot 15 \cdot 14! \\
&\equiv (-1) \cdot (-2) \cdot 14! \\
&\equiv 2 \cdot 14! \pmod{17}.
\end{align*}
Jadi
\[ 2 \cdot 14! \equiv -1 \pmod{17}. \]
Karena $2^{-1} \equiv 9 \pmod{17}$ (sebab $2 \cdot 9 = 18 \equiv 1$):
\[ 14! \equiv -1 \cdot 9 \equiv -9 \equiv 8 \pmod{17}. \]
Sisanya $8$.
\end{proof}

\begin{theorem}[Sisa Tiongkok / CRT]
Misalkan $m_1, \dots, m_k$ bilangan asli lebih dari $1$ yang saling relatif prima.
Maka untuk setiap bilangan bulat $a_1, \dots, a_k$, sistem kekongruenan
$x \equiv a_i \pmod{m_i}$ memiliki solusi yang tunggal modulo $m_1 m_2 \cdots m_k$.
\end{theorem}

\begin{example}
Tentukan bilangan asli terkecil $x$ dengan $x \equiv 2 \pmod{3}$ dan $x \equiv 3 \pmod{5}$.
\end{example}
\begin{proof}[Solusi]
Dari syarat pertama, $x = 3k + 2$ untuk suatu bilangan cacah $k$.
Substitusikan ke syarat kedua:
\begin{align*}
3k + 2 &\equiv 3 \pmod{5}, \\
3k &\equiv 1 \pmod{5}.
\end{align*}
Karena $3 \cdot 2 = 6 \equiv 1 \pmod{5}$, nilai $k$ terkecil adalah $k = 2$:
\[ x = 3 \cdot 2 + 2 = 8. \]
Periksa: $8 \equiv 2 \pmod{3}$ dan $8 \equiv 3 \pmod{5}$.
\end{proof}

\begin{example}[AIME 1983]
Carilah sisa pembagian $6^{83} + 8^{83}$ oleh $49$.
\end{example}
\begin{proof}[Solusi]
Tulis $6 = 7 - 1$ dan uraikan dengan binomial.
Semua suku dengan $7^2$ atau lebih tinggi habis dibagi $49$, jadi
\begin{align*}
6^{83} = (7-1)^{83} &\equiv (-1)^{83} + 83 \cdot 7 \cdot (-1)^{82} \\
&\equiv -1 + 581 \\
&\equiv 41 \pmod{49}.
\end{align*}
Dengan cara yang sama:
\begin{align*}
8^{83} = (7+1)^{83} &\equiv 1 + 83 \cdot 7 \\
&\equiv 43 \pmod{49}.
\end{align*}
Jumlahkan:
\[ 41 + 43 = 84 \equiv 35 \pmod{49}. \]
\end{proof}

\section{Bilangan Prima}

\begin{definition}
Bilangan asli $p > 1$ disebut \textbf{prima} apabila satu-satunya pembagi
positifnya adalah $1$ dan $p$ sendiri. Bilangan asli lebih dari $1$ yang bukan
prima disebut \textbf{komposit}.
\end{definition}

\begin{theorem}[Faktorisasi prima]
Setiap bilangan asli $n > 1$ dapat dinyatakan secara unik sebagai perkalian
bilangan prima:
\[ n = p_1^{\alpha_1} p_2^{\alpha_2} \cdots p_k^{\alpha_k} \]
dengan $p_1, \dots, p_k$ prima berbeda dan $\alpha_i$ bilangan asli.
\end{theorem}

\begin{example}
Tentukan banyak faktor positif dari $2004$.
\end{example}
\begin{proof}[Solusi]
Faktorkan dahulu:
\[ 2004 = 4 \cdot 501 = 2^2 \cdot 3 \cdot 167. \]
Setiap faktor positif berbentuk $2^a 3^b 167^c$ dengan
$0 \le a \le 2$, $0 \le b \le 1$, $0 \le c \le 1$.
Banyak pilihannya:
\[ (2 + 1)(1 + 1)(1 + 1) = 3 \cdot 2 \cdot 2 = 12. \]
\end{proof}

\begin{fact}[Sifat yang sering dipakai]
\begin{enumerate}
    \item Banyaknya bilangan prima tak hingga.
    \item Setiap bilangan prima lebih dari $3$ berbentuk $6k \pm 1$ untuk suatu bilangan asli $k$.
    \item Jika $p$ prima dan $p \mid ab$, maka $p \mid a$ atau $p \mid b$ (atau keduanya).
\end{enumerate}
\end{fact}

\begin{example}
Buktikan bahwa $24 \mid p^2 - 1$ untuk setiap bilangan prima $p > 3$.
\end{example}
\begin{proof}[Solusi]
Karena $p$ ganjil, $p - 1$ dan $p + 1$ dua bilangan genap berurutan,
jadi salah satunya habis dibagi $4$:
\[ 8 \mid (p-1)(p+1) = p^2 - 1. \]
Di sisi lain $p \equiv \pm 1 \pmod{3}$, sehingga
\begin{align*}
p^2 &\equiv 1 \pmod{3}, \\
3 &\mid p^2 - 1.
\end{align*}
Karena $8$ dan $3$ relatif prima:
\[ 24 \mid p^2 - 1. \]
\end{proof}

\section{Residu Pangkat}

Tidak semua sisa bisa muncul sebagai nilai suatu pangkat.
Keterbatasan ini sering menjadi kunci untuk menunjukkan suatu persamaan tidak punya solusi.

\begin{fact}[Residu kuadrat]
\begin{enumerate}
    \item Sisa $x^2$ dibagi $3$ hanya bisa $0$ atau $1$.
    \item Sisa $x^2$ dibagi $4$ hanya bisa $0$ atau $1$.
    \item Sisa $x^2$ dibagi $5$ hanya bisa $0$, $1$, atau $4$.
    \item Sisa $x^2$ dibagi $8$ hanya bisa $0$, $1$, atau $4$.
\end{enumerate}
\end{fact}

\begin{example}
Buktikan tidak ada bilangan asli $a$, $b$, $c$ yang memenuhi $a^2 + b^2 - 8c = 6$.
\end{example}
\begin{proof}[Solusi]
Tinjau persamaan itu modulo $8$:
\[ a^2 + b^2 \equiv 6 \pmod{8}. \]
Sisa kuadrat modulo $8$ hanya $0$, $1$, atau $4$,
sehingga jumlah dua sisa hanya bisa
\[ 0,\ 1,\ 2,\ 4,\ \text{atau}\ 5 \pmod{8}. \]
Nilai $6$ tidak pernah muncul, jadi persamaan itu tidak punya solusi.
\end{proof}

\begin{fact}[Residu pangkat lebih tinggi]
\begin{enumerate}
    \item Sisa $x^3$ dibagi $7$ hanya bisa $0$, $1$, atau $6$.
    \item Sisa $x^3$ dibagi $9$ hanya bisa $0$, $1$, atau $8$.
    \item Sisa $x^4$ dibagi $16$ hanya bisa $0$ atau $1$.
\end{enumerate}
\end{fact}

\begin{example}
Buktikan tidak ada bilangan bulat $x$ sehingga $x^3 \equiv 2 \pmod{7}$.
\end{example}
\begin{proof}[Solusi]
Hitung semua kemungkinan $x^3$ modulo $7$:
\begin{align*}
0^3 &\equiv 0 \pmod{7}, \\
1^3 &\equiv 1 \pmod{7}, \\
2^3 = 8 &\equiv 1 \pmod{7}, \\
3^3 = 27 &\equiv 6 \pmod{7}, \\
4^3 = 64 &\equiv 1 \pmod{7}, \\
5^3 = 125 &\equiv 6 \pmod{7}, \\
6^3 = 216 &\equiv 6 \pmod{7}.
\end{align*}
Sisa yang muncul hanya $0$, $1$, dan $6$.
Nilai $2$ tidak pernah muncul.
\end{proof}

\begin{example}
Carilah semua bilangan cacah $n$ sehingga $n^2 + 100$ merupakan bilangan kuadrat.
\end{example}
\begin{proof}[Solusi]
Misalkan $n^2 + 100 = m^2$ untuk suatu bilangan cacah $m$.
Pindahkan ruas dan faktorkan:
\[ (m - n)(m + n) = 100, \qquad m + n \ge m - n \ge 0. \]
Pasangan faktor $(m-n, m+n)$ yang jumlahnya genap hanyalah
\[ (2, 50) \quad \text{dan} \quad (10, 10), \]
yang memberi
\[ (m, n) = (26, 24) \quad \text{dan} \quad (10, 0). \]
Jadi $n = 0$ atau $n = 24$.
\end{proof}

\part*{Soal Latihan}

\section{Keterbagian}
\begin{enumerate}

	% Source: OSK SMA 2004 No. 4 (Pilihan Ganda)
	% Jawaban: B
	\item (OSK 2004)
	Untuk $a$ dan $b$ bilangan bulat dengan $a \ne 0$, notasi $a \mid b$ menyatakan
	``$a$ membagi $b$''. Pernyataan berikut yang \textbf{salah} adalah \dots
	\begin{enumerate}
		\item Jika $a \mid b$ dan $a \mid c$, maka $a \mid (bc)$.
		\item Jika $a \mid c$ dan $b \mid c$, maka $(ab) \mid c$.
		\item Jika $a \mid b$ dan $a \mid c$, maka $a \mid (b + c)$.
		\item Untuk setiap bilangan bulat $a \ne 0$ berlaku $a \mid 0$.
		\item Jika $a \mid b$, maka $a \mid (bc)$ untuk setiap bilangan bulat $c$.
	\end{enumerate}

	% Source: OSK SMA 2003 No. 1 (Pilihan Ganda)
	% Jawaban: 0
	\item (OSK 2003)
	Ada berapa banyak di antara bilangan-bilangan $20000002$, $20011002$,
	$20022002$, $20033002$ yang habis dibagi $9$?

	% Source: AIME 1984 (diadaptasi dari Modul Pra-OSK Keterbagian)
	% Jawaban: 592
	\item (AIME 1984)
	Bilangan asli $n$ adalah bilangan asli terkecil yang habis dibagi $15$
	dan semua digitnya sama dengan $8$ atau $0$. Tentukan nilai dari $\frac{n}{15}$.

	% Source: OSK SMA 2008 No. 5 (Pilihan Ganda)
	% Jawaban: 2
	\item (OSK 2008)
	Jumlah empat bilangan asli berturutan senantiasa habis dibagi $p$.
	Maka nilai $p$ terbesar adalah \dots

	% Source: OSP SMA 2015 No. 2 (Isian Singkat)
	% Jawaban: 4
	\item (OSP 2015)
	Banyaknya bilangan bulat $n$ sehingga $n + 1$ merupakan faktor dari $n^2 + 1$ adalah \dots

	% Source: OSK SMA 2013 No. 16 (Isian Singkat)
	% Jawaban: 64
	\item (OSK 2013)
	Diberikan himpunan $S = \left\{x \in \mathbb{Z} \;\middle|\; \frac{x^2 - 2x + 7}{2x - 1} \in \mathbb{Z}\right\}$.
	Banyaknya himpunan bagian dari $S$ adalah \dots

	% Source: OSK SMA 2014 No. 11 (Isian Singkat)
	% Jawaban: 595
	\item (OSK 2014)
	Bilangan-bilangan $1111$, $5276$, $8251$, dan $9441$ bersisa sama jika dibagi $N$.
	Nilai $N$ terbesar yang memiliki sifat tersebut adalah \dots

	% Source: OSK SMA 2012 No. 3 (Isian Singkat)
	% Jawaban: 960
	\item (OSK 2012)
	Bilangan terbesar $x$ kurang dari $1000$ sehingga terdapat tepat dua bilangan asli $n$
	dengan $\frac{n^2 + x}{n + 1}$ merupakan bilangan asli adalah \dots

	% Source: Finlandia MO 2001 (diadaptasi dari Modul Pra-OSK Keterbagian)
	% Jawaban: n = 6, 9, 2001
	\item (Finlandia MO 2001)
	Tentukan semua bilangan asli $n$ sedemikian hingga $n^2 + 2$ habis membagi $2001n + 2$.

	% Source: OSP SMA 2023 No. 2 (Isian Singkat)
	% Jawaban: 74
	\item (OSP 2023)
	Tentukan banyaknya bilangan asli $n \le 221$ sehingga
	\[ \frac{1^2 + 2^2 + \dots + n^2}{1 + 2 + \dots + n} \]
	merupakan bilangan bulat.

	% Source: Afrika Selatan MO 2021 (Modul Pra-OSK Keterbagian)
	% Jawaban: terkecil 17171, terbesar 98989
	\item (Afrika Selatan MO 2021)
	Tentukan bilangan asli terkecil dan terbesar dengan representasi digit
	yang berbentuk $ababa$ ($a \ne 0$) dan habis dibagi $11$.

	% Source: Kualifikasi Chile MO 2017 (Modul Pra-OSK Keterbagian)
	% Jawaban: d = 5 atau d = 8
	\item (Chile MO 2017)
	Dengan menggunakan dua digit berbeda $2$ dan $d$,
	bilangan $90$ digit $m = 2d22d222d2222d\ldots$ (blok berikutnya terdiri dari
	lima digit $2$ diikuti $d$, dan seterusnya) habis dibagi oleh $9$.
	Tentukan semua nilai yang mungkin dari $d$.

	% Source: Modul Pra-OSK Keterbagian
	% Jawaban: n = 2, 5, 8, 10, 15, 20, 26, 35, 40, 50, 65, 80, 90, 140, 170, 215, 290, 440, 890
	\item (Pra-OSK Keterbagian)
	Carilah semua bilangan asli $n$ sehingga $n + 10 \mid n^3 + 100$.

	% Source: India PRMO 2014 (Modul Pra-OSK Keterbagian)
	% Jawaban: 1
	\item (India PRMO 2014)
	Banyak bilangan asli $n$ di antara $1$ sampai $2014$ sehingga
	$\frac{8n}{9999 - n}$ merupakan bilangan bulat adalah \dots

	% Source: Modul Pra-OSK Keterbagian
	% Jawaban: n = k^2, k^2 + k, atau k^2 + 2k untuk suatu bilangan asli k
	\item (Pra-OSK Keterbagian)
	Tentukan semua bilangan asli $n$ sedemikian hingga $\lfloor \sqrt{n} \rfloor$ membagi $n$.

	% Source: Modul Pra-OSK Keterbagian
	% Jawaban: 1152
	\item (Pra-OSK Keterbagian)
	Diberikan sebuah barisan sepuluh bilangan asli naik tegas sedemikian hingga
	setiap bilangan membagi bilangan berikutnya.
	Jika jumlah kesepuluh bilangan tersebut adalah $1918$,
	maka semua kemungkinan bilangan kesepuluh dari barisan tersebut adalah \dots

	% Source: El Salvador MO 2020 (Modul Pra-OSK Keterbagian, notasi dirapikan)
	% Jawaban: 505
	\item (El Salvador MO 2020)
	Perhatikan bilangan $x_n = \overline{ab\,aab\,aaab\,\cdots\,\underbrace{a\ldots a}_{n}b}$,
	yaitu blok ke-$k$ terdiri dari $k$ digit $a$ diikuti satu digit $b$,
	dengan $a \ne 0$ dan $b$ keduanya digit.
	Tentukan banyak bilangan asli $n$ dengan $1 \le n \le 2020$ sedemikian hingga
	$x_n$ habis dibagi oleh $11$ untuk semua nilai $a$ dan $b$.

\end{enumerate}

\section{FPB dan KPK}
\begin{enumerate}[resume]

	% Source: Afrika Selatan MO 2006 (diadaptasi dari Modul Pra-OSK Keterbagian)
	% Jawaban: 70/37
	\item (Afrika Selatan MO 2006)
	Sederhanakan pecahan $\dfrac{2121212121210}{1121212121211}$ ke bentuk paling sederhana.

	% Source: OSK SMA 2024 No. 4 (Isian Singkat)
	% Jawaban: 214
	\item (OSK 2024)
	Misalkan $a$ dan $b$ bilangan bulat positif yang tidak memiliki faktor persekutuan selain $1$.
	Jika berlaku
	\[ \frac{1 + 2 + 3 + \dots + 104}{3 + 4 + 5 + \dots + 106} = \frac{a}{b}, \]
	maka nilai dari $a + b$ adalah \dots

	% Source: OSP SMA 2026 No. 5 (Isian Singkat)
	% Jawaban: 9
	\item (OSP 2026)
	Banyaknya pasangan terurut bilangan asli $(a, b)$ sehingga $a$ membagi $b^2$,
	$b$ membagi $a^2$, dan $ab = 10^{10}$ adalah \dots

	% Source: Modul Pra-OSK Keterbagian
	% Jawaban: 1 dan 12167
	\item (Pra-OSK Keterbagian)
	Hitunglah $FPB(133, 198)$ dan $FPB(6436343, 12167)$.

	% Source: Modul Pra-OSK Keterbagian
	% Jawaban: 32768
	\item (Pra-OSK Keterbagian)
	Berapa banyak pasangan bilangan asli $(x, y)$ yang memenuhi
	$FPB(x, y) = 5!$ dan $KPK(x, y) = 50!$?

	% Source: Modul Pra-OSK Keterbagian
	% Jawaban: 179
	\item (Pra-OSK Keterbagian)
	Definisikan $f(n) = FPB(n! + 1, (n + 1)!)$.
	Tentukan nilai dari $f(1) + f(2) + f(3) + \dots + f(30)$.

	% Source: OSK SMA 2021 Kemampuan Lanjut (Modul Pra-OSK Keterbagian)
	% Jawaban: 91
	\item (OSK 2021 Lanjut)
	Bilangan asli ganjil $b$ terbesar sedemikian hingga barisan bilangan asli
	$a_n = n^2 + 19n + b$ memenuhi
	\[ FPB(a_n, a_{n + 1}) = FPB(a_{n + 2}, a_{n + 1}) \]
	untuk setiap bilangan asli $n$ adalah \dots

	% Source: Modul Pra-OSK Keterbagian
	\item (Pra-OSK Keterbagian)
	Misalkan $m$ dan $n$ adalah bilangan asli yang memenuhi
	\[ FPB(m, n) + KPK(m, n) = m + n. \]
	Buktikan bahwa salah satu dari kedua bilangan itu membagi bilangan yang lainnya.

	% Source: Modul Pra-OSK Keterbagian
	\item (Pra-OSK Keterbagian)
	Diberikan bilangan asli $a$ dan $b$ yang memenuhi
	\[ \frac{a + 1}{b} + \frac{b + 1}{a} \]
	adalah bilangan bulat. Tunjukkan bahwa $FPB(a, b) \le \sqrt{a + b}$.

	% Source: Modul Pra-OSK Keterbagian
	\item (Pra-OSK Keterbagian)
	Misalkan $x$, $y$, $z$ adalah bilangan asli yang memenuhi
	\[ \frac{1}{x} - \frac{1}{y} = \frac{1}{z}. \]
	Misalkan $h = FPB(x, y, z)$.
	Tunjukkan bahwa $hxyz$ dan $h(y - x)$ keduanya merupakan bilangan kuadrat.

	% Source: Modul Pra-OSK Keterbagian
	% Jawaban: (t, t, t) dan semua permutasi dari (t, 2t, 3t) untuk t bilangan asli
	\item (Pra-OSK Keterbagian)
	Carilah semua tripel bilangan asli $(a, b, c)$ sedemikian hingga
	$a^3 + b^3 + c^3$ habis dibagi $a^2b$, $b^2c$, dan $c^2a$.

\end{enumerate}

\section{Modulo}
\begin{enumerate}[resume]

	% Source: OSK SMA 2009 No. 18 (Isian Singkat)
	% Jawaban: 6
	\item (OSK 2009)
	Jika $10^{999999999}$ dibagi oleh $7$, maka sisanya adalah \dots

	% Source: OSK SMA 2005 No. 8 (Pilihan Ganda, disajikan sebagai isian)
	% Jawaban: 6
	\item (OSK 2005)
	Tentukan digit terakhir dari $8^{8^{8^8}}$.

	% Source: OSK SMA 2003 No. 9 (Pilihan Ganda)
	% Jawaban: 13
	\item (OSK 2003)
	Misalkan $N$ adalah bilangan bulat terkecil yang bersifat bersisa $2$ jika dibagi $5$,
	bersisa $3$ jika dibagi $7$, dan bersisa $4$ jika dibagi $9$.
	Berapakah hasil penjumlahan digit-digit dari $N$?

	% Source: OSK SMA 2021 No. 3 (Isian Singkat)
	% Jawaban: 21
	\item (OSK 2021)
	Diberikan bilangan bulat positif $A$ dan $B$ yang bila dibagi $5$ berturut-turut bersisa $2$ dan $3$.
	Sisa pembagian $A(A + 1) + 5B$ jika dibagi $25$ adalah \dots

	% Source: OSK SMA 2023 No. 8 (Isian Singkat)
	% Jawaban: 30
	\item (OSK 2023)
	Sisa pembagian dari $5^{2021} + 11^{2022}$ oleh $64$ adalah \dots

	% Source: OSP SMA 2002 No. 16 (Isian Singkat)
	% Jawaban: 7
	\item (OSP 2002)
	Berapakah sisa pembagian $43^{43^{43}}$ oleh $100$?

	% Source: OSP SMA 2011 No. 17 (Isian Singkat)
	% Jawaban: semua bilangan asli
	\item (OSP 2011)
	Bilangan asli $n$ yang memenuhi $(-2004)^n - 1900^n + 25^n - 121^n$ habis dibagi $2000$ adalah \dots

	% Source: Canadian MO Qualification Repechage 2016 (diadaptasi dari Modul Pra-OSK Modulo)
	% Jawaban: tidak ada
	\item (Canadian MO 2016)
	Tentukan semua bilangan asli $n$ sehingga $11 \mid (3^n + 4^n)$.

	% Source: OSK SMA 2021 No. 3 (Isian Singkat)
	% Jawaban: 17
	\item (OSK 2021)
	Diketahui dua digit terakhir dari $a^{777}$ adalah $77$.
	Maka dua digit terakhir dari $a$ adalah \dots

	% Source: Modul Pra-OSK Modulo
	% Jawaban: 43
	\item (Pra-OSK Modulo)
	Tentukan dua digit terakhir dari $7^{707}$.

	% Source: Modul Pra-OSK Modulo
	% Jawaban: 9 dan 1
	\item (Pra-OSK Modulo)
	Carilah digit terakhir dari $13^{14}$ dan $13^{14^{15}}$.

	% Source: Modul Pra-OSK Modulo
	% Jawaban: 625
	\item (Pra-OSK Modulo)
	Tentukan sisa pembagian saat $1 \times 3 \times 5 \times \dots \times 2023$ dibagi oleh $1000$.

	% Source: Modul Pra-OSK Modulo
	% Jawaban: 76 dan 77
	\item (Pra-OSK Modulo)
	Carilah dua digit terakhir dari $1032^{1032}$ dan $17^{17^{17}}$.

	% Source: Argentina MO 2021 Level 1 (Modul Pra-OSK Modulo)
	% Jawaban: 9
	\item (Argentina MO 2021)
	Mikha menulis barisan aritmetika yang dimulai dari $1$ dengan selisih $3$, yakni
	\[ 1, 4, 7, 10, 13, 16, \ldots \]
	Kemudian, Fakhru menandai semua bilangan yang ditulis Mikha yang lebih dari $10$
	dan kurang dari $100000$ dan seluruh digitnya sama.
	Tentukan banyak bilangan yang ditandai oleh Fakhru.

	% Source: Modul Pra-OSK Modulo
	\item (Pra-OSK Modulo)
	Misalkan $m$, $a$, $b$ adalah bilangan bulat dengan $a$ dan $b$ relatif prima dengan $m$.
	Bilangan asli $x$ dan $y$ memenuhi $a^x \equiv b^x \pmod{m}$ dan $a^y \equiv b^y \pmod{m}$.
	Tunjukkan bahwa $a^{FPB(x, y)} \equiv b^{FPB(x, y)} \pmod{m}$.

	% Source: Modul Pra-OSK Modulo (pernyataan diperbaiki: p ganjil)
	\item (Pra-OSK Modulo)
	Misalkan $p$ adalah bilangan prima ganjil.
	Tunjukkan bahwa terdapat bilangan asli $x$ sedemikian hingga $x^2 + 1$
	habis dibagi oleh $p$ jika dan hanya jika $p \equiv 1 \pmod{4}$.

	% Source: Modul Pra-OSK Modulo (pernyataan diperbaiki: p - 1 digit)
	\item (Pra-OSK Modulo)
	Misalkan $p \ge 7$ adalah bilangan prima.
	Tunjukkan bahwa bilangan $11\ldots11$ (dengan $p - 1$ digit $1$) habis dibagi oleh $p$.

	% Source: Modul Pra-OSK Modulo
	\item (Pra-OSK Modulo)
	Misalkan $p$ adalah bilangan prima yang berbentuk $3k + 2$ yang membagi $a^2 + ab + b^2$
	untuk suatu bilangan bulat $a$ dan $b$.
	Buktikan bahwa $a$ dan $b$ keduanya habis dibagi oleh $p$.

	% Source: Modul Pra-OSK Modulo
	\item (Pra-OSK Modulo)
	Misalkan $n > 1$ adalah bilangan ganjil. Tunjukkan bahwa $n$ tidak membagi $3^n + 1$.

	% Source: Modul Pra-OSK Modulo
	\item (Pra-OSK Modulo)
	Misalkan $a$ adalah bilangan ganjil.
	Buktikan bahwa $a^{2^n} + 2^{2^n}$ dan $a^{2^m} + 2^{2^m}$ relatif prima saat $m \ne n$.

	% Source: Modul Pra-OSK Modulo (pernyataan diperbaiki: n > 1)
	\item (Pra-OSK Modulo)
	Misalkan $n > 1$ adalah bilangan asli dan $p$ bilangan prima sedemikian hingga
	$n$ membagi $p - 1$ dan $p$ membagi $n^3 - 1$.
	Tunjukkan bahwa $4p - 3$ adalah bilangan kuadrat.

	% Source: Modul Pra-OSK Modulo
	\item (Pra-OSK Modulo)
	Jika $p$ adalah bilangan prima ganjil dan $a$, $b$ adalah dua bilangan asli yang relatif prima,
	maka tunjukkan bahwa $FPB\left(a + b, \dfrac{a^p + b^p}{a + b}\right) = 1$ atau $p$.

	% Source: Modul Pra-OSK Modulo
	% Jawaban: (p - 1)p^2/2
	\item (Pra-OSK Modulo)
	Misalkan $p$ adalah bilangan prima ganjil.
	Untuk setiap $i = 1, 2, \ldots, p - 1$, notasikan $r_i$ sebagai sisa pembagian
	ketika $i^p$ dibagi dengan $p^2$.
	Carilah nilai dari $r_1 + r_2 + \dots + r_{p - 1}$.

	% Source: Modul Pra-OSK Modulo
	\item (Pra-OSK Modulo)
	Buktikan bahwa untuk setiap bilangan genap positif $n$,
	$n^2 - 1$ membagi $2^{n!} - 1$.

	% Source: Modul Pra-OSK Modulo
	\item (Pra-OSK Modulo)
	Misalkan $a$ dan $b > 2$ adalah bilangan asli.
	Tunjukkan bahwa $2^a + 1$ tidak habis dibagi oleh $2^b - 1$.

	% Source: Modul Pra-OSK Modulo
	\item (Pra-OSK Modulo)
	Misalkan $n > 2$ adalah bilangan asli.
	Tunjukkan bahwa $n^{n^{n^n}} - n^{n^n}$ habis dibagi oleh $1989$.

	% Source: Modul Pra-OSK Modulo
	\item (Pra-OSK Modulo)
	Misalkan $a$ dan $b$ adalah dua bilangan asli sedemikian hingga $b^n + n$
	merupakan kelipatan dari $a^n + n$ untuk setiap bilangan asli $n$.
	Tunjukkan bahwa $a = b$.

	% Source: Modul Pra-OSK Modulo
	\item (Pra-OSK Modulo)
	Misalkan $x$, $a$, $b$ adalah bilangan asli.
	Tunjukkan bahwa $FPB(x^a - 1, x^b - 1) = x^{FPB(a, b)} - 1$.

	% Source: Modul Pra-OSK Modulo
	\item (Pra-OSK Modulo)
	Misalkan $a$ adalah bilangan asli dan $p$ bilangan prima.
	Tunjukkan bahwa terdapat suatu bilangan asli $m$ yang memenuhi
	\[ m^{m^m} \equiv a \pmod{p}. \]

\end{enumerate}

\section{Bilangan Prima dan Pangkat}
\begin{enumerate}[resume]

	% Source: OSK SMA 2011 (diadaptasi dari Modul Pra-OSK Modulo)
	% Jawaban: p = 3
	\item (OSK 2011)
	Tentukan semua bilangan asli $p$ sedemikian hingga $p$, $p + 8$, dan $p + 16$
	semuanya merupakan bilangan prima.

	% Source: Modul Pra-OSK Modulo
	% Jawaban: (a, b) = (1, 1)
	\item (Pra-OSK Modulo)
	Tentukan semua bilangan asli $a$ dan $b$ sedemikian hingga $a^4 + 4b^4$
	merupakan bilangan prima.

	% Source: Modul Pra-OSK Modulo
	% Jawaban: (m, n) = (1, 3)
	\item (Pra-OSK Modulo)
	Carilah semua bilangan asli $m$ dan $n$ sedemikian hingga $6^m + 2^n + 2$
	adalah bilangan kuadrat.

	% Source: Modul Pra-OSK Modulo
	% Jawaban: n = 2
	\item (Pra-OSK Modulo)
	Tentukan semua bilangan asli $n$ sedemikian hingga $3n - 4$, $4n - 5$,
	dan $5n - 3$ semuanya merupakan bilangan prima.

	% Source: Modul Pra-OSK Modulo
	\item (Pra-OSK Modulo)
	Buktikan bahwa setiap bilangan prima yang lebih besar dari $3$
	berbentuk $6k \pm 1$ untuk suatu bilangan asli $k$.

	% Source: Modul Pra-OSK Modulo (soal yang sama sudah dibahas sebagai contoh di Materi)
	\item (Pra-OSK Modulo)
	Buktikan bahwa untuk setiap bilangan prima $p > 3$ berlaku $24 \mid p^2 - 1$.

	% Source: Modul Pra-OSK Modulo
	% Jawaban: 0
	\item (Pra-OSK Modulo)
	Banyak bilangan di barisan $11, 111, 1111, 11111, \ldots, 11\ldots11$
	dengan bilangan terakhir terdiri dari $2022$ digit $1$,
	yang merupakan bilangan kuadrat adalah \dots

	% Source: Modul Pra-OSK Modulo
	% Jawaban: n = 5
	\item (Pra-OSK Modulo)
	Tentukan semua bilangan asli $n$ sehingga $n! + 5$ merupakan bilangan kubik.

	% Source: Argentina MO 2021 Level 1 (Modul Pra-OSK Modulo)
	% Jawaban: (p, q) = (7, 3), (13, 3), (37, 3), (97, 3)
	\item (Argentina MO 2021)
	Tentukan semua bilangan prima $p$ dan $q$ yang lebih dari $1$ dan kurang dari $100$
	sedemikian hingga lima bilangan berikut:
	\[ p + 6,\quad p + 10,\quad q + 4,\quad q + 10,\quad p + q + 1 \]
	semuanya merupakan bilangan prima.

	% Source: Modul Pra-OSK Modulo
	\item (Pra-OSK Modulo)
	Jika $p$ adalah bilangan prima ganjil dan
	$\dfrac{a}{b} = \dfrac{1}{1} + \dfrac{1}{2} + \dots + \dfrac{1}{p - 1}$
	dalam bentuk paling sederhana,
	buktikan bahwa $p \mid a$.

	% Source: Modul Pra-OSK Modulo (pernyataan diperbaiki: memuat kata prima)
	\item (Pra-OSK Modulo)
	Buktikan bahwa jika $a^n + 1$ merupakan bilangan prima dengan $a > 1$,
	maka $a$ genap dan $n = 2^k$ untuk suatu bilangan cacah $k$.
	(Bilangan prima berbentuk $2^{2^k} + 1$ disebut bilangan prima Fermat.)

	% Source: Modul Pra-OSK Modulo (pernyataan diperbaiki: memuat kata prima)
	\item (Pra-OSK Modulo)
	Buktikan bahwa jika $a^n - 1$ merupakan bilangan prima dengan $a > 1$,
	maka $a = 2$ dan $n$ prima.
	(Bilangan prima berbentuk $2^n - 1$ disebut bilangan prima Mersenne.)

	% Source: Modul Pra-OSK Modulo
	\item (Pra-OSK Modulo)
	Buktikan bahwa untuk setiap bilangan asli $n$,
	terdapat $n$ buah bilangan asli berurutan yang tidak ada satupun
	di antaranya yang merupakan bilangan berpangkat sempurna.

	% Source: Modul Pra-OSK Modulo (redaksi diperbaiki)
	% Jawaban: ya, selalu ada
	\item (Pra-OSK Modulo)
	Apakah untuk setiap himpunan bilangan asli $\{a_1, a_2, \ldots, a_k\}$
	terdapat suatu bilangan asli $b$ sedemikian hingga setiap anggota himpunan
	$\{ba_1, ba_2, \ldots, ba_k\}$ merupakan bilangan berpangkat sempurna?

	% Source: OSP SMA 2010 (Modul Pra-OSK Modulo)
	% Jawaban: p = 13
	\item (OSP 2010)
	Tentukan semua bilangan prima $p$ sehingga $p^2 + 7^3$ merupakan bilangan kubik.

\end{enumerate}

\section{Referensi}
\begin{itemize}
    \item Modul Pra-OSK: Keterbagian (Azzam dan Akeyla).
    \item Modul Pra-OSK: Modulo (Azzam dan Akeyla).
\end{itemize}

\end{document}
